\section{Scope and thesis statement}\label{sec:intro-scope}
The locked thesis question is: \emph{is an
on-device Foundry-Local SLM traffic controller that matches or beats
classical control on real London topology, while producing a provable
accident audit, possible at all? And at what model scale and
configuration?} This is a feasibility and scale-threshold study:
everything runs in SUMO on three real London topologies plus synthetic
grids, only Euston demand magnitude is measured, realisations are
synthetic throughout, and the comparators are classical, not trained
deep-RL, since the target is feasibility rather than a new state of the
art in delay.

In answering it I argue something narrower than feasibility. Inside a
deterministic shield, what fine-tuning reliably moves is
\emph{authorship}, not delay. Stock students already author 39 to 62 per
cent of logged decision ticks, a majority on four of the five cells; the
fine-tuned students author more, on every cell I tested, though the size
of that shift varies widely and I did not isolate its mechanism. Delay
behaved differently. It moved materially only on the one network I
calibrated to clear, and not on every network that clears: the corridor I
calibrated did not follow. That gap is the thesis. An accident audit
concerns the certified model only so far as the model authored the
decisions, so authorship is the property that has to hold.
Chapter~\ref{ch:results} tests both halves separately, against both
classical comparators. Chapter~\ref{ch:discussion} reports where each
half holds, and Chapter~\ref{ch:conclusion} says where it fails.

